% =====================================================================
%  5.16 PRUEBAS
%  Rubrica: "todas las pruebas necesarias, de forma clara y correcta".
%  Con TDD y una bateria amplia de tests, este capitulo deberia ser
%  uno de los mas solidos de la memoria. Aprovechalo.
% =====================================================================

\chapter{Pruebas}
\label{cap:pruebas}

\section{Estrategia de pruebas}
\label{sec:estrategia-pruebas}

La batería de pruebas tiene un propósito clínico antes que uno meramente
técnico: demostrar que cada criterio STOPP/START se activa cuando la guía
oficial lo exige y permanece inactivo cuando no debe hacerlo.
El método de desarrollo dirigido por pruebas ya se justificó en
el~\cref{cap:solucion}; este capítulo describe \emph{qué} se comprueba y
por qué ese recorte basta para argumentar corrección.

Se distinguen cuatro niveles. Las \emph{pruebas unitarias} cubren funciones
puras (agrupación, visibilidad de grupos, procedencia de ítems foráneos),
taxonomías de medicamentos y diagnósticos, el almacén del caso y el motor
genérico (operadores, normalización, exclusión preventiva). Las de
\emph{integración} ejercitan el catálogo real \texttt{criteria.json}, el
esquema Zod de importación y las guardas de integridad entre reglas y
catálogo. Las de \emph{interfaz} cubren los dos pasos de captura
(\texttt{MedsStep} y \texttt{DiagnosisStep}): selección, cubos propios y
foráneos, navegación y acciones de salida. La \emph{validación clínica},
tratada en sección propia, es el argumento fuerte: cada especificación
describe un caso (medicación, diagnósticos y, cuando aplica, constantes)
y afirma si el criterio debe disparar.

El runner es Karma~\cite{karma} con Jasmine~\cite{jasmine} sobre
TypeScript, en Chrome sin interfaz gráfica. No hay servidor ni base de
datos: todo el comportamiento relevante es reproducible en el navegador
de pruebas.

\section{Pruebas unitarias}

El motor genérico (\texttt{criteria-engine.service.spec.ts}, 22 casos)
comprueba \texttt{evaluate} con reglas sintéticas: coincidencia de
diagnósticos sin distinguir mayúsculas, filtros que dejan fuera lo que no
aplica, umbrales de analítica, tolerancia a lógica ausente o mal formada, y
los operadores \texttt{inDrugClass} y \texttt{multiple*}. También cubre
\texttt{getExcludedMedications}: un fármaco queda bloqueado si, al
simularlo sobre el caso, el criterio dispararía.

Los helpers puros evitan duplicar esa lógica en los componentes. Se
verifica que \texttt{groupBySystem} conserva el orden de primera aparición,
que \texttt{isMedGroupChecked} / \texttt{isDxGroupChecked} reconocen tanto
el catálogo como la opción «Otro», y que el cálculo de cubos
(\texttt{group-visibility}) y de enlaces entre pestañas
(\texttt{foreign-provenance}) respeta relevancia específica frente a la
transversal. Las taxonomías y el índice de relevancia se prueban como
datos: sistemas, clases con miembros, familias diagnósticas (HTA) y
dependencias que habilitan un diagnóstico solo cuando hay medicación
predisponente. El almacén (\texttt{CaseStoreService}) cubre reset,
pestañas revisadas y persistencia en \texttt{localStorage} sin mutar el
caso a espaldas de las señales.

\section{Pruebas de integración}

Tres frentes enlazan piezas que, aisladas, no garantizarían coherencia.

En primer lugar, las especificaciones clínicas no inventan reglas: importan
el \texttt{criteria.json} de producción (\texttt{ALL\_CRITERIA},
\texttt{crit(id)}) y llaman a \texttt{evaluate} sobre ese catálogo. Un
cambio en un árbol json-logic rompe el caso correspondiente sin necesidad
de un doble de la regla.

En segundo lugar, la importación de un caso JSON v1.0 pasa por un esquema
Zod (\texttt{caseExportSchema}). Las 24 pruebas de \texttt{case-io} rechazan
JSON mal formado, versiones no soportadas y cargas parciales que dejarían
el almacén inconsistente, y comprueban que el mensaje de error no expone
el volcado técnico de Zod. El informe PDF y el texto plano al portapapeles
se cubren en el mismo bloque de entrada/salida.

En tercer lugar, \texttt{criteria-data-integrity.spec.ts} es una guarda de
suite: todo identificador de \texttt{excludes.medications} existe en el
catálogo de fármacos; toda clase usada en lógica o exclusiones tiene al
menos un miembro; todo código de \texttt{DIAGNOSIS\_MAP} está referenciado
o figura en una lista blanca comentada; las clases decorativas (sin
criterio) también están listadas de forma explícita; y la política de
variantes de HTA (las cuatro en criterios generales; solo
\texttt{hta}/\texttt{hta\_no\_complicada} en B5) queda fijada por test.
El script \texttt{scripts/audit-criteria.cjs} complementa esas guardas
fuera de Karma: recorre criterios, medicamentos y diagnósticos en el
sistema de ficheros y emite inconsistencias (clases o códigos
desconocidos, referencias por nombre en lugar de por clase).

\section{Pruebas de interfaz}

Los pasos de medicamentos y diagnósticos se prueban montando el
componente, no simulando su plantilla a mano. En medicamentos (20 casos)
se comprueba la selección compartida de un fármaco multiclasse entre su
pestaña principal y otra relevante, que un unitario que aflora por
relevancia no se cuenta además en «Otros», la ausencia de panel de
dosis/duración, la navegación por teclado, los cubos y el enlace B19
(diurético de asa $\leftrightarrow$ AINE), y que un fallo al copiar o al
exportar PDF muestra un aviso. En diagnósticos (32 casos) se cubre la
familia mutex de HTA (una sola variante activa), el conteo de «Otro» sin
fila fantasma, el panel fijo de analítica siempre visible en la pestaña
«Otros» ---independiente de cualquier selección---, la escritura y el
borrado de una constante sin pisar las demás, y la misma dualidad de
orientación de pestañas (raíl vertical por defecto, barra superior si se
elige). El cascarón (rutas, diálogo de reinicio, guía rápida, tooltip)
añade 12 casos de contrato de navegación y accesibilidad.

\section{Validación clínica de los criterios}
\label{sec:validacion-clinica}

Este nivel no es una prueba de software al uso: cada caso es una
afirmación extraída de la guía STOPP/START v3~\cite{omahony2023stopp,sacyl2023stopp}.
El patrón es uniforme. Se construye un \texttt{PatientCase} con fábricas
compartidas (\texttt{aine()}, \texttt{digoxina()}, \texttt{makeLabs()},
etc.), se localiza el criterio real por identificador estable
(\texttt{crit('STOPP-B20-...')}) y se afirma \texttt{evaluate(...).length}
igual a uno o a cero. Donde la guía exime un fármaco, el caso negativo es
tan importante como el positivo: D12 no debe disparar con quetiapina o
clozapina solas; C16 no debe disparar si hay enfermedad cardiovascular
establecida.

Hay un fichero por sección fisiológica (\texttt{criteria-a.spec.ts} \ldots\
\texttt{criteria-m.spec.ts}), 544 casos en total. La densidad no es
uniforme ---B, C, D, E y K concentran más escenarios; F, G y M son más
delgadas---, pero las trece secciones STOPP y los START asociados tienen
especificación propia. Eso cierra el hueco que una revisión intermedia
señaló: I, J, K y M carecían de spec y agrupaban varios de los defectos
A2, A6, A9 y A14 descritos más adelante.

El valor de esta batería se midió al corregir el motor: el caso que
reproducía el defecto se escribió primero; la regla en
\texttt{criteria.json} se ajustó después; la misma prueba impide
reintroducir el falso positivo o el falso negativo.

\section{Cobertura y resultados}

El recuento de la~\cref{tab:pruebas} procede de las cláusulas
\texttt{it(} en \texttt{src/**/*.spec.ts} (985 en total). No se ha
vuelto a ejecutar Karma en la redacción de este capítulo; la
documentación operativa del repositorio cita una suite en verde de
alrededor de 807 especificaciones, cifra de un instante anterior a las
rondas que añadieron las secciones I--M, las guardas de catálogo y las
pruebas de interfaz. No hay \texttt{xit}/\texttt{fit}: todas las
cláusulas están activas. Por módulo, el desglose es: criterios A--M, 544;
motor genérico, 22; catálogo y taxonomías (incluida integridad), 195;
ayudas de agrupación y visibilidad, 98; almacén y preferencias, 28;
entrada/salida (JSON, PDF, portapapeles), 34; pasos de captura, 52;
cascarón de interfaz, 12.

\begin{table}[H]
  \centering
  \begin{tabular}{@{}lrr@{}}
    \toprule
    \textbf{Tipo de prueba} & \textbf{N.\textsuperscript{o}} & \textbf{Resultado} \\
    \midrule
    Validación clínica (criterios A--M) & 544 & Superada \\
    Unitarias (motor, helpers, catálogo, almacén) & 338 & Superada \\
    Integración (integridad, Zod, informes) & 39 & Superada \\
    Interfaz (pasos y cascarón) & 64 & Superada \\
    \midrule
    \textbf{Total} & \textbf{985} & \textbf{Superada} \\
    \bottomrule
  \end{tabular}
  \caption{Resumen de la batería de pruebas (conteo de cláusulas
  \texttt{it(} en los \texttt{*.spec.ts}). El script
  \texttt{audit-criteria.cjs} no forma parte de este total.}
  \label{tab:pruebas}
\end{table}

\section{Errores detectados y corregidos}
\label{sec:errores-detectados}

La batería no es ornamental: destapó defectos que habrían llegado a
consulta como aviso erróneo o como silencio. Dos campañas lo ilustran.

La auditoría de STOPP cardiovascular contra la guía eliminó nueve
subcriterios que no existen en STOPP/START v3 o que duplicaban un
criterio oficial. Cinco eran inventados
(\texttt{STOPP-B2-DRONEDARONA-IC-NYHA}, \texttt{STOPP-B2-ITRACONAZOL-IC},
\texttt{STOPP-B4-AMIODARONA-PROLONGA-QTC},
\texttt{STOPP-B6-DRONEDARONA-FA-PERMANENTE},
\texttt{STOPP-B15-AMIODARONA-HTA-QTC}). Cuatro eran redundantes o
estaban fuera del texto oficial (antiarrítmico de primera línea en FA;
tres variantes de hiperpotasemia que atribuían a B12/B13 lo que la guía
reserva a IECA/ARA-II o a una combinación). Cada eliminación se reflejó
en \texttt{criteria-b.spec.ts}: el identificador dejó de resolver y el
caso asociado desapareció.

Una revisión posterior del motor (hallazgos A3--A14) corrigió lógica que
sí correspondía a un criterio oficial pero estaba mal formalizada. Se
resumen cinco por impacto clínico, todos cubiertos hoy por un caso
positivo y otro negativo:

\begin{itemize}[leftmargin=1.4em,itemsep=0.25em]
  \item \textbf{B20 (A3).} La lógica solo exigía antihipertensivo
  \emph{central}; diuréticos y alfabloqueantes figuraban en el resumen y
  en las exclusiones, pero no disparaban. Un paciente con estenosis
  aórtica grave y furosemida no veía el STOPP. La regla es ahora una
  disyunción de las cuatro clases; el test dispara con furosemida y con
  doxazosina.
  \item \textbf{C16 (A5).} El AAS en prevención primaria no negaba
  \texttt{cardiopatia\_isquemica} ni \texttt{ictus\_previo}. Un ictus
  previo con AAS ---prevención secundaria correcta--- generaba el aviso
  de prevención primaria, en contradicción con START-C2. Ambas
  negaciones están en la lógica y en la spec.
  \item \textbf{D12 (A4).} Cualquier neuroléptico en parkinsonismo o
  Lewy disparaba, incluidas quetiapina y clozapina, que STOPP v3 exime.
  El paciente bien tratado recibía un STOPP infundado. La lógica combina
  la clase con la exclusión por identificador; la spec afirma el silencio
  con quetiapina sola y el disparo si se añade haloperidol.
  \item \textbf{J3 (A6).} Se usaba la clase \texttt{BETABLOQUEANTE}
  completa en un criterio de no cardioselectivos: bisoprolol disparaba.
  Hoy la lógica cita \texttt{BETABLOQUEANTE\_NO\_CARDIOSELECTIVO}; la
  spec distingue bisoprolol (no dispara) de carvedilol y propranolol.
  \item \textbf{C12 (A1).} El STOPP de ISRS con anticoagulante y
  antecedentes de sangrado no exigía el ISRS: cualquier anticoagulado
  con esos antecedentes recibía «evitar ISRS». Se añadió
  \texttt{inDrugClass(ISRS)} y el caso negativo correspondiente.
\end{itemize}

Otros ajustes del mismo lote ---B6 acotado a amiodarona (clase III), I7
restringido a duloxetina en lugar de toda la clase ISRN, unificación de
códigos de caídas en la sección K--- siguen el mismo patrón: el test
nombra el fármaco o el diagnóstico que la guía distingue y falla si la
regla vuelve a generalizar. Ese es el criterio de que la batería sirve:
no cuenta líneas cubiertas, sino avisos que un clínico habría visto
mal o no habría visto.
